\begin{center}
	{\Large \textbf{Appendix A}}-\vspace*{0.5 cm}
	{\Large \textbf{Sample Code}}
	
\end{center}
\begin{lstlisting}[language=Python, caption=AWS Scanner Core Logic]
class AWSScanner(CloudScanner):
    def scan_resources(self) -> List[CloudResource]:
        resources = []
        methods = [self._scan_ec2_instances, self._scan_lambda_functions,
                   self._scan_rds_instances, self._scan_s3_buckets,
                   self._scan_sqs_queues, self._scan_iam_roles]
        for method in methods:
            try:
                resources.extend(method())
            except Exception as e:
                logger.error(f"Scan error: {str(e)}")
        return resources

    def _get_ec2_dependencies(self, instance_id: str) -> Set[str]:
        dependencies = set()
        try:
            instance = self.ec2_client.describe_instances(InstanceIds=[instance_id])['Reservations'][0]['Instances'][0]
            if 'VpcId' in instance:
                dependencies.add(instance['VpcId'])
            for sg in instance.get('SecurityGroups', []):
                dependencies.add(sg['GroupId'])
            if 'SubnetId' in instance:
                dependencies.add(instance['SubnetId'])
        except Exception as e:
            logger.error(f"EC2 dep error: {str(e)}")
        return dependencies
\end{lstlisting}

\begin{lstlisting}[language=Python, caption=AI Remediation Integration]
def analyze_vulnerabilities_with_gemini(report_json: dict, gemini_api_key: str) -> dict:
    prompt = (
        "Analyze the following cloud security report JSON and generate a structured vulnerability analysis. "
        "Your output must be valid JSON with an array named 'vulnerabilities'. Each vulnerability object must contain the following keys: "
        "'priority' (either 'critical' or 'non critical'), 'problem', 'solution', 'cli_command', and 'file'. "
        "Input: " + json.dumps(report_json) + " Output:"
    )
    
    url = f"https://generativelanguage.googleapis.com/v1beta/models/gemini-1.5-flash:generateContent?key={gemini_api_key}"
    headers = {"Content-Type": "application/json"}
    data = {"contents": [{"parts": [{"text": prompt}]}]}
    
    response = requests.post(url, headers=headers, json=data)
    # ... parsing logic ...
    return structured_output
\end{lstlisting}
